\section{CT18 的输出：PDF、QCD 参数、部分子光度、矩}
\label{sec:OverviewCT18}
本节综述 CT18 PDF 及相应的部分子光度、Mellin 矩和 QCD 拉氏量参数的行为。鉴于图数量庞大，对 CT18Z 拟合，本节只展示最关键的比较；CT18Z PDF 的其余对照图在 Appendices~\ref{sec:ATL7ZWchi2} 与 \ref{sec:LMCT18Z} 中给出。

\subsection{作为 $x$ 与 $Q$ 函数的部分子分布}
\subsubsection{各味夸克的 PDF}


图 \ref{fig:ct18pdf} 给出 $Q = 2$ 与 $100$ GeV 处 CT18 部分子分布函数的概览。函数 $x f(x,Q)$ 作为 $x$ 的函数绘出，涵盖 $u,
\overline{u}, d, \overline{d}, s = \overline{s}$ 与 $g$ 各味。
我们假定 $s(x,Q_0)=\bar s(x,Q_0)$，因为其差异与零相容且不确定度很大~\cite{Lai:2007dq}。图中展示对表~\ref{tab:EXP_1} 与 \ref{tab:EXP_2} 所列全球数据的中心拟合，对应于我们所选 PDF 参数化下最低的总 $\chi^2$。误差带表示 90\% 置信水平（C.L.）的 PDF 不确定度。
\begin{figure}[p]
	\center
	\includegraphics[width=0.49\textwidth]{images/u_100_CT18.pdf}
	\includegraphics[width=0.49\textwidth]{images/d_100_CT18.pdf}
	\includegraphics[width=0.49\textwidth]{images/ubar_100_CT18.pdf}
	\includegraphics[width=0.49\textwidth]{images/dbar_100_CT18.pdf}
	\includegraphics[width=0.49\textwidth]{images/s_100_CT18.pdf}
	\includegraphics[width=0.49\textwidth]{images/g_100_CT18.pdf}
	\caption{$Q=100$ GeV 处 CT18（紫色实线）、\CTHERAII（灰色短划线）与 CT18Z（品红长划线）NNLO 组合 90\% C.L. PDF 不确定度的比较。不确定度带按中心 CT18 NNLO PDF 归一。
		\label{fig:PDFbands1}}
\end{figure}


从 CT14$_\mathrm{HERAII}$ NNLO 到 CT18 NNLO PDF 的相对变化，最直观的方式是比较它们各自的 PDF 不确定度。图~\ref{fig:PDFbands1} 比较了关键味在 90\% C.L. 的 PDF 误差带，每条带按相应的最优拟合 CT18 NNLO PDF（实线紫色曲线/带）归一。品红长划线与灰色短划线曲线/带分别是 $Q=100$ GeV 处的 CT18Z 与 CT14$_\mathrm{HERAII}$ NNLO PDF。图~\ref{fig:PDFbands2} 把同样的误差带各自归一到其中心拟合，以便直接比较它们的 PDF 不确定度。
\begin{figure}[p]
	\center
	\includegraphics[width=0.49\textwidth]{images/u_100_CT18_band.pdf}
	\includegraphics[width=0.49\textwidth]{images/d_100_CT18_band.pdf}
	\includegraphics[width=0.49\textwidth]{images/ubar_100_CT18_band.pdf}
	\includegraphics[width=0.49\textwidth]{images/dbar_100_CT18_band.pdf}
	\includegraphics[width=0.49\textwidth]{images/s_100_CT18_band.pdf}
	\includegraphics[width=0.49\textwidth]{images/g_100_CT18_band.pdf}
	\caption{同图~\ref{fig:PDFbands1}，但每个误差带按其各自的中心 PDF 归一，以便直接比较 CT18 NNLO 与 CT18Z、CT14$_\mathrm{HERAII}$ NNLO 的 PDF 不确定度。
		\label{fig:PDFbands2}}
\end{figure}

我们对 NNLO PDF 作出若干观察。与 \CTHERAII{} 相比，CT18 $u$ PDF 在几乎所有 $x$ 值处都略小，最大降幅出现在 $x \sim 10^{-3}$。$d$ PDF 在 $x \sim 10^{-3}$ 与 $x \sim 0.2$ 处增大，而在 $ x\sim 0.01$ 处略降。$\bar u$ 与 $\bar d$ 分布都在 $x \sim 0.3$ 处变小、在 $x \sim 0.05$ 处变大，只是 $\bar d$ 的降幅更大。此外，除 $x\sim 0.2$ 处的 $d$ PDF 外，$u$、$d$、$\bar u$、$\bar d$ 的误差带与 \CTHERAII{} 大致相同。
中心奇异性（$s$）PDF 在 $x < 0.01$ 增大、在 $0.2< x < 0.5$ 减小；后者区间内奇异夸克 PDF 在 CT18 中基本不受约束，正如 \CTHERAII~NNLO 一样。而且其不确定度带在 $x > 10^{-4}$ 处略大于 \CTHERAII，这是更灵活参数化以及纳入 LHC 数据的结果。我们已核查：驱动上述夸克与反夸克 PDF 变化的最重要数据集是 LHCb $W$ 与 $Z$ 玻色子数据，见表~\ref{tab:EXP_2} 中 Exp.~IDs=250、245 与 246，重要性依此排序。纳入 LHCb $W,Z$ 数据之后，再加 CMS 8 TeV $W$ 电荷不对称数据（Exp.~ID=249）仅使 CT18 PDF 发生极轻微变化。中心胶子 PDF 在 CT18 中于 $x\approx 0.3$ 处减小，且 $x \sim 0.1$ 及以下误差带更窄。$g$ PDF 在 $0.1 < x < 0.4$ 的下降源于 CMS 与 ATLAS 喷注数据（Exp.~IDs=545、543 与 544，按此顺序）以及 ATLAS 8 TeV $Z$ 玻色子横向动量（$p_T$）数据（Exp.~ID=253）的纳入。在已含 LHC 喷注数据集的情况下，再把 ATLAS 与 CMS 顶夸克对数据（Exp.~IDs=580 与 573）加入拟合不会带来统计上显著的 PDF 变化。

\subsubsection{PDF 的比值 \label{sec:PDFratios}}
\begin{figure}[t]
	\center
	\includegraphics[width=0.49\textwidth]{images/du_1p4_CT18.pdf}
	\includegraphics[width=0.49\textwidth]{images/du_100_CT18.pdf}
	\includegraphics[width=0.49\textwidth]{images/pdfs_CT18NNLO_CT14HERA2NNLO_CJ15nlo__10_0GeV_A90CL__10___dou__pdf__lin-lin_ect.pdf}
	\includegraphics[width=0.49\textwidth]{images/pdfs_CT18NLO_CT14HERA2NLO_CJ15nlo__10_0GeV_A90CL__10___dou__pdf__lin-lin_ect.pdf}

	\caption{上排：$Q=1.4$ 与 $100$ GeV 处 CT18、\CTHERAII、CT18Z NNLO 组合比值 $d(x,Q)/u(x,Q)$ 的 90\% C.L. 不确定度。下排：同样内容，比较 $Q=10$ GeV 处 CT18 与 \CTHERAII{} NNLO 比值（左下）及各自 NLO 比值（右下）与 CJ15 NLO 组合。
		\label{fig:DOUband}}
\end{figure}

现在考察各种 PDF 的比值，从图s.~\ref{fig:DOUband} 与 \ref{fig:DOUband2} 所示的 $d/u$ 开始。与 \CTHERAII{} 相比，CT18 中 $d/u$ 的变化可概括为：中心比值在 $x > 0.5$（$x<10^{-2}$）处降低（升高），且 $x < 10^{-2}$ 处不确定度减小。超过 $x=0.5$ 后 $d/u$ 的误差带增大；而自 CT14 NNLO~\cite{Dulat:2015mca} 起采用的参数化形式保证 $d/u$ 在 $x\rightarrow 1$ 时趋于常数，这与众多核子结构的理论模型预言一致。这通过令 $u_v$ 与 $d_v$ PDF 的 $(1-x)^{a_2}$ 指数相等实现，即 $a^{u_v}_2 = a^{d_v}_2$（见 App.~\ref{sec:AppendixParam}）。该选择只影响超出拟合数据覆盖范围的极高 $x$（$x\gtrsim 0.9$）外推。在 $x<0.9$，我们的参数化足够灵活，即使没有这一约束也能覆盖解并复现拟合的不确定度带。例如图~\ref{fig:DOUband} 的不确定度带在 $x=0.9$ 处延伸到 $d/u=0$。在可及的 $x$ 范围内，它也覆盖了我们具有独立 $a^{u_v}_2$ 与 $a^{d_v}_2$ 的候补最优拟合。若无此选择，单个被拟合 Hessian PDF 组的 PDF {\it 比值}将不具备外推到 $x = 1$ 处有限常数的参数自由度；相反，哪怕拟合参数的微小差异 $a^{u_v}_2 \neq a^{d_v}_2$ 都会使其在最高 $x$ 点发散或趋于 0，产生 $d/u$ 上的无穷大不确定度——这与该区域的电子-强子经验数据及常见强子结构模型都不相容。在不约束 $a^{u_v}_2$ 与 $a^{d_v}_2$ 相等的拟合中，二者也可能因输入参数的数值精度而碰巧相等，但据我们大量经验，这种巧合几乎从不发生。

类似的逻辑也适用于其他 PDF 比值，包括下面描述的海夸克分布的低 $x$ 形式。如前所述，CT18 中 $u$、$d$、$\bar u$、$\bar d$ 夸克的参数化形式与 \CTHERAII{} 相同。

\begin{figure}[t]
	\center
	\includegraphics[width=0.49\textwidth]{images/dou_1p4_CT18_band.pdf}
	\includegraphics[width=0.49\textwidth]{images/dou_100_CT18_band.pdf}
	\caption{类似图~\ref{fig:DOUband} 上排关于 $d/u$ 的面板，但把每个拟合并按其各自中心值归一以比较 PDF 不确定度。
		\label{fig:DOUband2}}
\end{figure}

在如此高的 $x$ 处，CTEQ-JLab 分析（CJ15）~\cite{Accardi:2016qay} 通过纳入 CT18 因选择割断 $W > 3.5 \mbox{ GeV}$ 而排除的较低 $W$、较高 $x$ 固定靶 DIS 数据，并考虑在该运动学区域重要的高阶扭度与核效应，独立确定了 NLO 的 $d/u$ 比值。图~\ref{fig:DOUband} 显示 CT18 的中心预言在 $x > 0.1$ 处不同于 CJ15。CT 与 CJ 的不确定度带彼此相符，尽管 CJ15 的误差带远小于 CT18——这部分归因于 CJ15 采用 $\Delta \chi^2 = 1$ 判据。由于 CJ15 PDF 只提供 $\alpha_s$ 的 NLO 版本，我们在图~\ref{fig:DOUband} 左下（右下）框中把 CJ15 NLO $d/u$ 比值与相应的 CT18 NNLO（NLO）比值比较。

\begin{figure}[tb]
	\center
	\includegraphics[width=0.49\textwidth]{images/dbub_1p4_CT18.pdf}
	\includegraphics[width=0.49\textwidth]{images/dbub_100_CT18.pdf}
	\includegraphics[width=0.49\textwidth]{images/Rs_1p4_CT18.pdf}
	\includegraphics[width=0.49\textwidth]{images/Rs_100_CT18.pdf}
	\caption{$Q=1.4$ 或 $100$ GeV 处比值 $\bar d(x,Q)/\bar u(x,Q)$ 与 $\left(s(x,Q)+\bar
		s(x,Q)\right)/\left(\bar u(x,Q) +\bar d(x,Q)\right)$ 的 90\% C.L. 不确定度比较，涉及 CT18（蓝色实线）、CT18Z（品红长划线）与 \CTHERAII~NNLO（灰色短划线）组合。
		\label{fig:DBandSBbands}}
\end{figure}


\begin{figure}[tb]
	\center
	\includegraphics[width=0.49\textwidth]{images/dbub_1p4_CT18_err.pdf}
	\includegraphics[width=0.49\textwidth]{images/dbub_100_CT18_err.pdf}
	\includegraphics[width=0.49\textwidth]{images/Rs_1p4_CT18_err.pdf}
	\includegraphics[width=0.49\textwidth]{images/Rs_100_CT18_err.pdf}
	\caption{同图~\ref{fig:DBandSBbands}，但按各自中心拟合归一，给出 $Q=1.4$ 或 $100$ GeV 处比值 $\bar d(x,Q)/\bar u(x,Q)$ 与 $\left(s(x,Q)+\bar
		s(x,Q)\right)/\left(\bar u(x,Q) +\bar d(x,Q)\right)$ 的 90\% C.L. 不确定度（CT18 蓝色实线、CT18Z 品红长划线、\CTHERAII~NNLO 灰色短划线）。
		\label{fig:DBandSBbands2}}
\end{figure}

转向图~\ref{fig:DBandSBbands} 中海夸克 PDF 的比值，我们观察到左插图中小 $x$ 处 $\bar
d(x,Q)/\bar u(x,Q)$ 的不确定度在 CT18 中有所减小。对于 $x\! >\! 0.1$，CT18 对 $\bar u$ 与 $\bar d$ 的非微扰参数化形式保证比值
$\bar d(x,Q_0)/\bar u(x,Q_0)$ 可趋于常数，且中心拟合中该常数接近 1。$\bar d/\bar u$ 的不确定度也已减小，最显著的是 $x\! \gtrsim\! 10^{-3}$ 区域，主要由于 LHCb 数据集（Exp.~IDs=250、245 与 246）的纳入，参见 $Q\!=\!1.4\,\mathrm{GeV}$ 处图~\ref{fig:DBandSBbands2} 左上面板。

在高 $Q$ 值处（如图~\ref{fig:DBandSBbands} 中 $Q=100$ GeV 的右侧面板），这些比值既依赖夸克 PDF 本身，也同样依赖 $Q_0$ 处大 $x$ 胶子的行为。因此，对于在极大 $x$ 和 $Q_0$ 处胶子 PDF 增强、海夸克 PDF 受抑的 CT18Z，比值 $\bar d/\bar u$ 与 $R_s$ 的不确定度在 $x\gtrsim 0.8$ 与大 $Q$ 处反而减小，反映出在此 $\{x,Q\}$ 区域主要驱动海夸克 PDF 的 $g\rightarrow q\bar q$ 劈裂的味对称性。

$x < 0.03$ 处奇异性 PDF 的总体增大以及 $x  < 10 ^{-3}$ 处 $\bar u$、$\bar d$ PDF 的减小（参见图~\ref{fig:PDFbands1}）导致奇异-非奇异海夸克 PDF 之比更大：
\begin{equation}
	R_s(x,Q) \equiv \frac{s(x,Q) + \bar{s}(x,Q)}{\bar{u}(x,Q) + \bar{d}(x,Q)}\ ,
\label{eq:Rs}
\end{equation}
见图~\ref{fig:DBandSBbands}。
$R_s(x,Q)$ 度量味-$\mathrm{SU}(3)$ 对称破缺随 $x$ 和 $\Q$ 的依赖性；较早的分析通常固定 $R_s = 0.5$。较近的一些 CTEQ 研究~\cite{Lai:2007dq,Olness:2003wz} 考察了当时对 $R_s$ 的约束，主要由 CCFR 与 NuTeV 合作组的中微子诱导 SIDIS 双缪子产生测量驱动，也包括精确的单举 HERA 测量。这些工作发现了 $s^+(x)\! \equiv\! s(x)\! +\! \bar{s}(x)$ 具有区别于 $\bar{u}\!+\!\bar{d}$ 的独立 $x$ 依赖性的显著证据，但无法排除趋于零的奇异性动量分数不对称
$\langle x \rangle_{s^-}\! =\! \int^1_0 dx\ x [s-\bar{s}](x,\Q=m_c)\! =\! 0$。

在本工作中，我们继续假定 $s^-(x,Q)=0$，聚焦 $s^+(x,Q)$ 及相关的 $R_s(x,Q)$；这两个量共同反映较早的带电流 DIS 数据与新 LHC 测量之间的相互作用，将在 Sec.~\ref{sec:Quality} 与 App.~\ref{sec:AppendixCT18Z} 详细讨论。这里提一句：在 $x  \ll 10 ^{-3}$ 处，$R_s$
比值完全由参数化形式决定；CT10 发现它与 PDF 味的精确 $\mathrm{SU}(3)$ 对称一致，即在 $x\rightarrow 0$ 时 $R_s(x,Q) \rightarrow 1$，但不确定度很大。
$x\rightarrow 0$ 的 $\mathrm{SU}(3)$ 对称渐近解在 CT14 与 \CTHERAII{} 中均未被强制，因此其 $R_s$ 比值在 $x\approx 10^{-5}$、$Q=1.4$ GeV 处约为 0.3 到 0.5。
在 CT18 中，我们采用了不同的 $s$-PDF 非微扰参数化形式（增加一个参数），但仍保证 $R_s$ 在 $x\! \to\! 0$ 的稳定行为：CT18 与 CT18Z 拟合中 $R_s (x \to 0)$ 分别约为 0.7 与 1。

\subsubsection{PDF 的 $x$ 依赖性变化小结}
可将各具体过程对中心 CT18 拟合的拉动总结如下。

\begin{itemize}

\item LHC 单举喷注产生对中心胶子 PDF $g(x,Q)$ 最显著的总体影响是在原 PDF 不确定度带内于 $x > 0.2$ 处使其温和减小。喷注数据集的拉动在对部分系统误差去关联后（参见 Sec.~\ref{sec:DataJets}），以及在加入理论值的 $0.5\%$ MC 不确定度后，几乎不变。各喷注数据集对 $g(x,Q)$ 的拉动在整个 $x$ 范围内并不遵循统一趋势，各种测量之间也不一致，例如由图~\ref{fig:L2glu} 的 $L_2$
敏感度及 Sec.~\ref{sec:QualityOverview} 的 LM 扫描所示。

\item LHCb 数据对所有过程合并后对 $u$、$d$ 与 $s$ 夸克有一定影响，并在小 $x$ 处向上拉动 $s (x,Q)$。

\item ATLAS 8 TeV $Z$ $p_T$ 数据（Exp.~ID=253）在 CT18 NNLO 拟合假定的名义 QCD 标度下弱弱地把 $x>0.05$ 的胶子 PDF 向下拉，方向与 LHC 单举喷注数据的平均拉动相似。与喷注实验相比，这些数据拉动的相对大小可由图~\ref{fig:L2glu} 关于 $g(x,Q)$ 的 $L_2$ 敏感度图估计。

\item ATLAS 7 TeV $W$ 与 $Z$ 快度分布数据（Exp.~ID=248）只包含在 CT18A 与 CT18Z 中，对 PDF 影响最大，讨论见
App.~\ref{sec:AppendixCT18Z}。其拉动方向与 LHCb 类似。

\item LHC 的 $t \bar t$ 双微分截面数据似乎也偏好大 $x$ 处更软的胶子，但拉动不具统计显著性，即远弱于数据点数目多得多的单举喷注数据。

\end{itemize}

Sec.~\ref{sec:QualityOverview} 将综合运用多种统计手段进一步深入探讨这些约束。

\subsection{$\alpha_s$ 与 $m_c$ 的全球拟合}
\label{sec:AlphasDependence}

\begin{figure}[tbp]
\centering
\includegraphics[width=0.49\textwidth]{images/alphas_scanTct18_pjb05a.pdf} \ \
\includegraphics[width=0.46\textwidth]{images/alphas_scanTct18_pjb05a2.pdf}
\caption{NNLO 下 CT18 强耦合常数在 $M_Z$ 标度的扫描。左：全部数据集合在一起的 $\chi^2$ 变化（粗黑线）以及对 $\alpha_s(M_Z)$ 拉动力特别强的几个实验各自的 $\chi^2$ 变化。右：CT18 所拟合全部实验的 $\chi^2$ 变化值，分别归入 DIS、DY 与 top/jets 组合数据集。右面板``Jets+top''曲线的 $\Delta \chi^2$ 增长主要由喷注产生的约束驱动。虽然 $t\bar{t}$ 产生对 $\alpha_s$ 有重要敏感度，但顶夸克数据点数较少，使其在完整拟合中的影响居中。
	}
\label{fig:lm_alphas}
\end{figure}

{\bf QCD 耦合常数的确定}。按照长期沿用的做法~\cite{Lai:2010nw}，在 CT18 这类规范 PDF 组中，$\alpha_s(M_Z)$ 取为世界平均值 $\alpha_s(M_Z)\! =\! 0.118$~\cite{Tanabashi:2018oca}；同时对该中心值上下的一系列固定 $\alpha_s(M_Z)$ 生成替代 PDF（即 ``$\alpha_s$ 系列''），用以评估组合的 PDF+$\alpha_s$ 不确定度。
%
文献~\cite{Lai:2010nw} 展示了如何在全球拟合中评估组合的 $\mathrm{PDF} \, + \, \alpha_s$ 不确定度。
如那里所示，$\alpha_s$ 的变动一般会引起偏好 PDF 参数的补偿性调整（关联），以保持与那些同时约束 $\alpha_s$ 与 PDF 的实验数据集的符合。同时，可以定义一个量化所有关联效应的``$\alpha_s$ 不确定度''。随着全球 QCD 数据集规模增长，越来越多的实验通过硬截面的辐射修正或标度破坏（尤其是在很宽的物理标度范围 $\Q$ 上）引入对 $\alpha_s(M_Z)$ 的敏感性。

考察单个实验乃至整个实验组对 $\alpha_s$ 敏感度的最佳方式，也许就是观察当 $\alpha_s(M_Z)$ 取值变动时其 $\chi^2$ 的变化。CT18 NNLO 与 CT18 NLO 对 $\alpha_s(M_Z)$ 的扫描分别示于图s.~\ref{fig:lm_alphas} 与 \ref{fig:lm_alphas_nlo}。在本节与下一节所有展示扫描的图中，我们都绘出一族曲线
\begin{equation}
  \Delta \chi_E^2(a) \equiv \chi_E^2(a)-\chi_E^2(a_0),
  \label{DelChi2Scan}
\end{equation}
作为某个参数 $a$ 的函数。变动量 $\Delta \chi^2_E(a)$ 是实验 $E$ 在横轴所示固定 $a$ 值处的 $\chi^2$ 值（其中 $\chi^2_E(a)$ 已对其余自由参数边缘化）与 $a$ 取完整 CT18 数据集全局 $\chi^2$ 极小值 $a=a_0$ 时之差。图中对全部实验（记作``Total''或
``$\chi^2_{\rm tot}$''）以及所示 $a$ 范围内 $\Delta \chi^2_E$ 变动最大的前几名实验绘出 $\Delta \chi^2$ 曲线。于是按定义 $\Delta \chi^2_{\rm tot}(a_0) =0$。

需要指出，本次扫描中 $\alpha_s$ 在所有精确辐射贡献中被变动，但在制表的 $K$ 因子中保持固定。这一近似大大简化计算；我们已核实，在被拟合的 $\alpha_s$ 范围内，它引起的 $\chi^2$ 变化仅为高阶不确定度的很小一部分。

\begin{figure}[tbp]
\centering
\includegraphics[width=0.8\textwidth]{images/i2Tn2as0_11872_DEchi2_re2_ect.pdf}
\caption{
	同图~\ref{fig:lm_alphas}，但现在显示 NLO 精度下 $\alpha_s(M_Z)$ 的扫描。
        }
\label{fig:lm_alphas_nlo}
\end{figure}

从图~\ref{fig:lm_alphas} 可见，各数据集对 $\alpha_s(M_Z)$ 的中心值及其不确定度有不同的敏感度。根据扫描结果，对 $\alpha_s(M_Z)$ 敏感度最高的是 HERA I+II 数据集，其次是 BCDMS 质子数据。相对于完整 CT18 数据集，这两个实验都偏好较低的 $\alpha_s(M_Z)$，约在 $0.114-0.116$，但不确定度更宽。这两个 DIS 数据集对 $\alpha_s(M_Z)$ 的依赖主要通过标度破坏效应，但其庞大的数据点数量与实验和理论的精度带来了很大的敏感度。

LHC 单举喷注产生（尤其 CMS 7 与 8 TeV 数据）一般偏好较大的 $\alpha_s(M_Z)$，ATLAS 8 TeV $Z$ $p_T$ 数据也是如此。完整的 CT18 数据集偏好 $\alpha_s(M_Z, \mbox{NNLO})=0.1164 \pm 0.0026$（68\% C.L.），采用``全局容忍度''方案定义，对应 $\Delta \chi^2\! =\! 37$ 区间（相应 90\% 区间由 $\Delta \chi^2\! =\! 100$ 定义）。用 CT18Z 提取的 $\alpha_s(M_Z)$ 十分相近：$0.1169\pm0.0026$，参见图~\ref{fig:lm_alphasz}。这些数值可与 CT14 用较小的 HERA+LHC 数据集得到的 $\alpha_s(M_Z)\!=\!0.1150^{+0.0036}_{-0.0024}$ 比较。

完整数据集的 $\Delta\chi^2$ 分布非常接近抛物线形，单个数据集则差些。数据集合集（例如全部 DIS 数据、全部 DY 数据、全部喷注与顶夸克数据）的 $\Delta\chi^2$ 曲线也呈抛物线形，符合中心极限定理的预期。从图~\ref{fig:lm_alphas} 右面板可以清楚看到，DIS 数据整体偏好比 DY 对产生、喷注与顶夸克产生更小的 $\alpha_s(M_Z)$。
因此 $\alpha_s$ 不确定度的确切大小并未被很好地确定，依赖于约定，因为各（类）实验之间的拉动在几十个 $\chi^2$ 单位的水平上并不一致。

扫描练习也可以在 $\alpha_s$ 的 NLO 下进行，如图~\ref{fig:lm_alphas_nlo} 所示。事实上，NLO 与 NNLO 结果之差可以作为对其确定的理论不确定度的一个部分估计。尽管不确定度与 NNLO 相近，中心值略高：$\alpha_s(M_Z,\mbox{NLO})=0.1187\pm 0.0027$。我们注意到，NLO 下对 $\alpha_s(M_Z)$ 具有领先敏感度的实验之间的定性相互关系与 NNLO 大体相同：HERA 合并数据（Exp.~ID=160）与 BCDMS $F^p_2$ 数据（Exp.~ID=101）再次偏好较低值，而 ATLAS 7 TeV 喷注数据（Exp.~ID=544）与 8 TeV
$Z$ $p_T$ 数据（Exp.~ID=253）向相反方向拉动，且 NLO 比 NNLO 更强。NLO 下 $\alpha_s$ 值高出约 0.002 的偏好与其他 PDF 组的发现一致~\cite{Ball:2011us,Harland-Lang:2015nxa,Abramowicz:2015mha,Alekhin:2017kpj}。

总之，我们发现 CT18 数据集偏好的 $\alpha_{s}(M_Z)$ 比 CT14 更大，名义不确定度略小。


\begin{figure}[tbp]
\centering
\includegraphics[width=0.6\textwidth]{images/i2Tn3Q01-00mc1-30_DEchi2_re_ect.pdf}
\caption{
	用 CT18 数据集对粲夸克极点质量 $m_c$ 数值在 NNLO 作 $\chi^2$ 扫描。拟合设定见正文。此 $m_c$ 扫描对应的 CT18Z 版本见 App.~\ref{sec:AppendixCT18Z} 图~\ref{fig:lm_mcZ}。
        }
\label{fig:lm_mc}
\end{figure}

{\bf 约束粲极点质量}。
对微扰理论的其他输入也可开展类似研究，例如粲夸克极点质量 $m_c$。对粲质量依赖性的决定性研究超出了本文的范围：实验对 $m_c$ 的偏好可能受初始标度 $Q_0$、重夸克方案的辅助设定以及非微扰粲的可能性影响~\cite{Gao:2013wwa,Hou:2017khm,Dulat:2013hea}。在我们所做的候补拟合中观察到，传统取值 $m_c^{pole}=1.3$ GeV 仍与 CT18(Z) 全球数据相容，但最近的 HERA DIS 与 LHC 矢量玻色子产生实验总体上可能轻度偏好 1.4 GeV 或更高的极点质量。

所观察到的趋势可在图~\ref{fig:lm_mc} 中看到：我们展示了在不同极点质量 $m_c$ 下对 CT18 数据集作 NNLO 拟合时总数据集与主导实验的 $\chi^2$ 变化。为把 $m_c$ 依赖性与 $Q_0$ 依赖性分开，我们取 $Q_0=1$ GeV，并使用更灵活的（在小 $x$ 处符号不定）胶子参数化，以更好容纳如此低 $Q$ 下解的全部范围。取 $Q_0=1$ GeV（CT10 的 $m_c$ 依赖性研究~\cite{Gao:2013wwa} 也这样做）使我们能扩大所考察的 $m_c$ 范围，而 $Q=1$ GeV 处胶子的额外灵活性则是容纳 $Q>1.3$ GeV 处 CT18 PDF 不确定度全范围所必需的。从粗黑曲线的极小可见，扫描偏好 $m_c\! =\! 1.3$ GeV；该质量比仅由 HERA 合并粲产生数据（Exp.~ID=147）单独给出的偏好略大，否则后者暗示
$m_c\!\gtrsim\!1.2$ GeV。另一方面，HERA Run I 与 II 合并单举数据基本给出 $m_c$ 的下限，并偏好更大的量级 $m_c\! > \!1.45$ GeV。
不过这些偏好相当弱：在大范围的 $m_c$ 上总 $\chi^2$ 变化只有十个单位。图~\ref{fig:lm_mc} 中其他实验的单个敏感度还要更弱。
应当指出，纳入 ATLAS 7 TeV $W/Z$ 数据（Exp.~ID=248）及其他与 CT18Z 相关的改变会使图~\ref{fig:lm_mc} 的图像重新配置并提高 $m_c$ 的最优拟合值，见图~\ref{fig:lm_mcZ} 及 App.~\ref{sec:AppendixCT18Z} 的讨论。
同理，若取 $Q_0=m_c$（另一个可接受的选择），拉动的模式也会略有不同。


\subsection{LHC 的部分子光度
\label{sec:PDFLuminosities}
}

图~\ref{fig:lumia} 展示用 CT18 与 CT18Z NNLO PDF 计算的 LHC 14 TeV 部分子光度，并与先前的 \CTHERAII{} 版本对比。为了只在物理可及区域内比较光度，我们在计算光度积分时限制末态绝对快度在 5 个单位以内，参见 \cite{Hou:2016sho} 式 (28)。在每个味组合的比较中，我们再次给出两种归一化结果：左侧图按共同参考（CT18 NNLO 或 NLO）归一，或按各自的中心预言归一。

与单个 PDF 的情形一样，部分子光度的 CT18 中心结果仍与 \CTHERAII{} 非常接近。另一方面，各 PDF 组合的光度 PDF 不确定度有一定缩减，尤其是含胶子的光度。在希格斯玻色子质量区 $M_X\! \sim\! 125$ GeV，图~\ref{fig:lumia} 下面板所示的 $gg$ 光度 $L_\mathit{gg}$ 改进很小。
但在 TeV 尺度的质量范围内，$L_\mathit{gg}$ PDF 不确定度的缩减更为可观，接近 $\sim\!20\%$。
用 CT18Z NNLO 计算的部分子光度在若干方面明显不同于 CT18。例如在
$W/Z$ 玻色子质量区，CT18Z 的 $qq$ 光度中心预言比 CT18 高约 $3\!-\!4\%$。其他部分子光度在低质量区 $M_X\! \lesssim\! 100$ GeV 也同样在 CT18Z 中增强，主要缘于 CT18Z 使用的依赖 $x$ 的 DIS 因子化标度。这种小 $x$ 增强在 CT18X 与 Z 中大致相同，而 CT18 更接近 \CTHERAII。
高质量区的夸克-夸克光度 $L_{qq}$ 在 CT18Z 中相对未变，
而 $L_{gq}$ 与 $L_{gg}$ 在 $M_X\! \gtrsim\! 100\!-\!300$ GeV 受抑制；对胶子-胶子光度，在 100 GeV 到 1 TeV 之间抑制可达 $\sim\!4\%$。

在图s.~\ref{fig:lumiCT18NLOvsothers}
与~\ref{fig:lumiCT18NNLOvsothers} 中，我们将这些部分子光度与其他组比较：
CJ15~\cite{Accardi:2016qay}、MMHT14~\cite{Harland-Lang:2014zoa} 与
NNPDF3.1~\cite{Ball:2017nwa}。此处我们对 CT18 PDF 采用重标定的 68\% C.L.，以匹配其他组的约定。
图~\ref{fig:lumiCT18NLOvsothers} 的比较在 NLO 进行，
因为 CJ15 PDF 没有 QCD NNLO 版本。
CJ15 NLO 光度的 PDF 不确定度小于 CT18 NLO（见右插图），部分原因在于其更小的容忍判据（$\Delta \chi^2 =1$）以及 $\bar u$、$\bar d$ PDF 采用的参数化形式欠灵活。
在质量 $10\lesssim M_{X}\lesssim 2\cdot  10^{3}$ GeV 范围，CT18 NLO $qq$ 光度中心值比 CJ15 高约 8\% 到 5%；更高质量时低 8--10\%（见左插图）。
在质量范围 $20\lesssim M_{X}\lesssim 2\cdot  10^{3}$ GeV 内，CT18 的 $L_{qq}$ PDF 误差带覆盖 CJ15 的误差带。
在与希格斯产生相关的质量区域
$100\lesssim M_{X}\lesssim 300$ GeV，CT18 NLO $gq$ 光度比 CJ15 高约 2%。在其他区域也更高：低质量 $M_X\lesssim 100$ GeV 处高 8%，高质量 $M_X \gtrsim 650$ GeV 处差异约 12%（在 $M_X \approx 5$ TeV）。对 $gg$ 光度，CT18 NLO 处处更高，特别地 $M_X \approx 5$ TeV 处差异大于 20%。

用 CT18 NNLO PDF 得到的光度与 MMHT14、NNPDF3.1 PDF 得到者在图~\ref{fig:lumiCT18NNLOvsothers} 比较。
NNPDF3.1 PDF 选定 $\alpha_s(M_Z) = 0.118$。
三个组的 $qq$ 与 $gq$ 光度中心值在质量区 $100\lesssim
M_{X}\lesssim 10^{3}$ GeV 内彼此相差百分之几，且 PDF 不确定度相当。
比较质量范围 $20 < M_X < 300$ GeV 的 NNLO $gg$ 光度，我们看到 MMHT14 与 CT18 差别在百分之一左右，而 NNPDF3.1 高 $2-3\%$。此外，MMHT14 在此范围内的不确定度带与 CT18 相似，NNPDF3.1 则较小。
在更大质量 $M_X \gtrsim 300$ GeV，我们观察到 NNPDF3.1 光度迅速下降。而且 NNPDF3.1 的不确定度在全质量范围都更小。

2012 至 2015 年间，针对当时的全球 PDF 拟合开展了若干细致研究~\cite{Ball:2012wy, Rojo:2015acz, Butterworth:2015oua}，包括方法学的基准比对。由此获得的理解促成了 2015 年关于 LHC 上 PDF 使用的建议~\cite{Butterworth:2015oua}。基准比对也改善了 CT、MMHT 与 NNPDF PDF 之间的一致性，进而使 PDF4LHC 工作组得以把这些全球 PDF 组合起来，形成广泛使用的 PDF4LHC15 PDF 组合。

自 2015 年以来，大量 LHC 数据进入了最新的全球拟合。
在某些情形中，这使得各组中心 PDF 之间的差异比 2015 年 PDF 发布版更大。这一变化可能归因于多种因素。特别地，小 $x$ 重求和或 NNLO DIS 截面中非传统的 QCD 标度选择都会修改小 $x$ PDF，图s.~\ref{fig:PDFbands1} 与 \ref{fig:lumia} 中的 CT18Z NNLO 即为例证。NNPDF3.1 的部分子光度在某些区域与 CT18、MMHT2014 的差别更加明显，参见图~\ref{fig:lumiCT18NNLOvsothers}。一项后续研究正在进行，以更好地理解 LHC 数据与方法学选择对各全球 PDF 的影响。





\begin{figure}[!htbp]
	\begin{center}
		\hspace*{-0.4cm}\includegraphics[width=0.48\textwidth]{images/qq_CT18.pdf}
		\includegraphics[width=0.48\textwidth]{images/qq_CT18_err.pdf} \\
		\includegraphics[width=0.48\textwidth]{images/gq_CT18.pdf}
		\includegraphics[width=0.48\textwidth]{images/gq_CT18_err.pdf} \\
		\includegraphics[width=0.48\textwidth]{images/gg_CT18.pdf}
		\includegraphics[width=0.48\textwidth]{images/gg_CT18_err.pdf}
	\end{center}
	\vspace{-2ex}
	\caption{
		$\sqrt{s} = 14$ TeV 下中央快度区 $|y|<5$ 内各过程的部分子光度：$L_{qq}$（上面板）、$L_{gq}$（中间面板）
		与 $L_{gg}$（下面板）；采用 CT18（紫色实线）、CT18Z（灰色短划线）
		与 CT14$_\mathrm{HERAII}$（品红长划线）NNLO PDF 计算。左面板给按 CT18 归一的光度比，右面板给出每种光度的误差带（每个 PDF 组合按其自身中心预言归一）。
	}
\label{fig:lumia}
\end{figure}



\begin{figure}[!htbp]
	\begin{center}
		\hspace*{-0.4cm}
		\includegraphics[width=0.48\textwidth]{images/Lumi_14TeV_ym-lt-5-0_asym_20__68CL-CT18NLO_CJ15nlo__00_qQr_ect.pdf}
		\includegraphics[width=0.48\textwidth]{images/Lumi_14TeV_ym-lt-5-0_asym_20__68CL-CT18NLO_CJ15nlo__00_qQ2_ect.pdf} \\
		\includegraphics[width=0.48\textwidth]{images/Lumi_14TeV_ym-lt-5-0_asym_20__68CL-CT18NLO_CJ15nlo__00_gqr_ect.pdf}
		\includegraphics[width=0.48\textwidth]{images/Lumi_14TeV_ym-lt-5-0_asym_20__68CL-CT18NLO_CJ15nlo__00_gq2_ect.pdf} \\
		\includegraphics[width=0.48\textwidth]{images/Lumi_14TeV_ym-lt-5-0_asym_20__68CL-CT18NLO_CJ15nlo__00_ggr_ect.pdf}
		\includegraphics[width=0.48\textwidth]{images/Lumi_14TeV_ym-lt-5-0_asym_20__68CL-CT18NLO_CJ15nlo__00_gg2_ect.pdf}

	\end{center}
	\vspace{-2ex}
	\caption{同图~\ref{fig:lumia}，比较 CT18 NLO 与 CJ15 NLO 的部分子光度。
	}
\label{fig:lumiCT18NLOvsothers}
\end{figure}

\clearpage

\begin{figure}[!htbp]
	\begin{center}
		\hspace*{-0.4cm}
		\includegraphics[width=0.48\textwidth]{images/Lumi_14TeV_ym-lt-5-0_asym_30__68CL-CT18NNLO_MMHT2014nnlo_NNPDF31nnloas01181000__00_qQr_ect.pdf}
		\includegraphics[width=0.48\textwidth]{images/Lumi_14TeV_ym-lt-5-0_asym_30__68CL-CT18NNLO_MMHT2014nnlo_NNPDF31nnloas01181000__00_qQ2_ect.pdf} \\
		\includegraphics[width=0.48\textwidth]{images/Lumi_14TeV_ym-lt-5-0_asym_30__68CL-CT18NNLO_MMHT2014nnlo_NNPDF31nnloas01181000__00_gqr_ect.pdf}
		\includegraphics[width=0.48\textwidth]{images/Lumi_14TeV_ym-lt-5-0_asym_30__68CL-CT18NNLO_MMHT2014nnlo_NNPDF31nnloas01181000__00_gq2_ect.pdf} \\
		\includegraphics[width=0.48\textwidth]{images/Lumi_14TeV_ym-lt-5-0_asym_30__68CL-CT18NNLO_MMHT2014nnlo_NNPDF31nnloas01181000__00_ggr_ect.pdf}
		\includegraphics[width=0.48\textwidth]{images/Lumi_14TeV_ym-lt-5-0_asym_30__68CL-CT18NNLO_MMHT2014nnlo_NNPDF31nnloas01181000__00_gg2_ect.pdf}

	\end{center}
	\vspace{-2ex}
	\caption{同图~\ref{fig:lumia}，比较 $\alpha_s(M_Z)=0.118$ 的 CT18、MMHT'2014 与 NNPDF3.1 NNLO 部分光度。
	}
\label{fig:lumiCT18NNLOvsothers}
\end{figure}


\clearpage


\subsection{PDF 矩与求和规则}
\label{sec:moments}

积分 PDF {\it Mellin} 矩的知识长期以来备受关注，既因其唯象学用途，也因其与强子结构的格点 QCD 计算相关~\cite{Lin:2017snn,Hobbs:2019gob,Lin:2020rut}。
就前者而言，PDF 矩可作为比较各种全球分析与理论途径的有价值基准，也是描述 PDF 本身的信息性指标。这尤其是因为给定阶数的 PDF 所得的数值结果与底层部分子分布的 $x$ 依赖性相连：一般而言，更高阶的 Mellin 矩主要由 PDF 的高 $x$ 行为决定。文献~\cite{Hobbs:2019gob} 分析了 HEP 数据对格点可计算量——特别是 Mellin 矩与部分子准分布函数——的敏感度，以进一步发展方兴未艾的 PDF-格点事业~\cite{Lin:2017snn,Lin:2020rut}。

只要所论矩在其全支撑范围内收敛，原则上任何唯象 PDF 的积分矩都可以由其底层分布算得。但在本分析中，我们特别关注

\begin{equation}
\langle x^n\rangle_g(Q)=\int_{0}^{1}dx\ x^{n}\,  g(x,\Q) \, ,
\label{eq:gmom}
\end{equation}
（胶子取 $n\!=\!1$），以及夸克分布的
\begin{align}
	\langle x^n\rangle_{q^+}(Q) &= \int_{0}^{1}dx\ x^{n}\,  [q + \overline{q}](x,\Q) \ \ \mathrm{for}\ \ n = 1, 3, \dots \nonumber \\
	\langle x^n\rangle_{q^-}(Q) &= \int_{0}^{1}dx\ x^{n}\,  [q - \overline{q}](x,\Q) \ \ \mathrm{for}\ \ n = 2, 4, \dots
\label{eq:qmom}
\end{align}
其中电荷共轭偶（奇）夸克组合记作 $q^\pm = q \pm \overline{q}$。
我们主要考虑式~(\ref{eq:gmom}) 中 $n\!=\!1$ 的特定 PDF 矩以及式~(\ref{eq:qmom})，以便与格点 QCD 的确定兼容——后者只能计算 $q\!+\!\bar{q}$ 型分布的奇（$n\! =\! 1, 3, \dots$）矩与 $q\!-\!\bar{q}$ 型分布的偶（$n\! =\! 2, 4, \dots$）矩。
这是因为格点计算从强子矩阵元提取积分 Mellin 矩：
\begin{equation}
	\frac{1}{2}\sum_{s}\langle p,s|\mathcal{O}_{\{\mu_{1},\cdots,\mu_{n+1}\}}^{q}|p,s\rangle = 2 \langle x^{n+1} \rangle_{q}\,[p_{\mu_{1}}\cdots p_{\mu_{n+1}}-{\rm traces}]\ .
\label{eq:OPE}
\end{equation}
式~(\ref{eq:OPE}) 中的格点算符是
$\mathcal{O}_{\{\mu_{1},\cdots,\mu_{n+1}\}}\! \sim\! i^n \overline{q} \gamma_{\mu_1} \overleftrightarrow{D}_{\mu_2} \cdots \overleftrightarrow{D}_{\mu_{n+1}} q$，其中的协变导数使得逐次插入导数既提高提取矩的阶数，又交替改变其在电荷共轭下的偶奇性。


\begin{table}
\begin{tabular*}{\textwidth}{c| @{\extracolsep{\fill}} cc|c|ccc}
\hline
PDF moment                             &   {\bf CT18}    &    {\bf CT18Z}   &    {\bf \CTHERAII}   &  {\bf MMHT14}   &   {\bf CJ15}      &     {\bf NNPDF3.1}    \tabularnewline
\hline
$\langle x \rangle_{u^+-d^+}$          &  $0.156(7)$     &    $0.156(6)$    &    $0.159(6)$     &   $0.151(4)$    &    $0.1518(13)$   &   $0.152(3)$          \tabularnewline
$\langle x^2 \rangle_{u^--d^-}$        &  $0.055(2)$     &    $0.055(2)$    &    $0.055(2)$     &   $0.053(2)$    &    $0.0548(2)$    &   $0.057(3)$          \tabularnewline
$\langle x^3 \rangle_{u^+-d^+}$        &  $0.022(1)$     &    $0.022(1)$    &    $0.022(1)$     &   $0.022(1)$    &    $0.0229(1)$    &   $0.022(1)$          \tabularnewline
\hline
$\langle x \rangle_{g}$                &  $0.414(8)$     &    $0.407(8)$    &    $0.415(8)$     &   $0.411(9)$    &    $0.4162(8)$    &   $0.410(4)$          \tabularnewline
\hline
$\langle x \rangle_{u^+}$              &  $0.350(5)$     &    $0.350(4)$    &    $0.351(5)$     &   $0.348(5)$    &    $0.3480(6)$    &   $0.348(4)$          \tabularnewline
$\langle x \rangle_{d^+}$              &  $0.193(5)$     &    $0.194(5)$    &    $0.193(6)$     &   $0.197(5)$    &    $0.1962(9)$    &   $0.196(4)$          \tabularnewline
$\langle x \rangle_{s^+}$              &  $0.033(9)$     &    $0.041(8)$    &    $0.031(8)$     &   $0.035(8)$    &    $0.0313(2)$    &   $0.039(4)$          \tabularnewline
\hline
$\langle x^2 \rangle_{u^-}$            &  $0.085(1)$     &    $0.084(1)$    &    $0.085(1)$     &   $0.083(1)$    &    $0.0853(2)$    &   $0.085(3)$          \tabularnewline
$\langle x^2 \rangle_{d^-}$            &  $0.030(1)$     &    $0.029(1)$    &    $0.030(1)$     &   $0.030(1)$    &    $0.0305(2)$    &   $0.028(3)$          \tabularnewline
$\langle x^2 \rangle_{s^-}$            &    ---          &      ---         &      ---          &   $0.001(1)$    &      ---          &   $0.001(4)$          \tabularnewline
\hline
$\langle 1 \rangle_{\bar{d}-\bar{u}}$  &  $-0.12(35)$    &    $-0.07(29)$   &    $-0.37(41)$    &   $0.084(15)$   &    $0.103(20)$    &     ---               \tabularnewline
$\kappa_s$                             &  $0.49(16)$     &    $0.61(14)$    &    $0.46(13)$     &   $0.51(14)$    &      ---          &   $0.563(82)$         \tabularnewline
\end{tabular*}\caption{
	我们汇集了按 CT18、CT18Z、CT14$_\mathrm{HERAII}$、MMHT14、CJ15 NLO 与
	NNPDF3.1 计算的若干 PDF 矩数值，全部在标度 $Q=2$ GeV。选取这些矩既因其可作为唯象计算的基准，也因其与格点 QCD 计算相关。自上而下依次为：同位矢量（$u\!-\!d$）分布的三个最低阶矩、胶子的一阶矩、分味 $u, d, s$ 分布的第一、二阶矩，以及轻夸克味对称破缺的两个度量：味 $\mathrm{SU}(2)$ 差的零阶矩 $\langle 1 \rangle_{\bar{d}-\bar{u}}$
	和与奇异性压低相关的矩比 $\kappa_s$（定义见式~(\ref{eq:str-supp})）。注意 NNPDF3.1 未约束 $\bar{d}$ 与
	$\bar{u}$ 分布在 $x\!\to\!0$ 处重合，故 $\langle 1 \rangle_{\bar{d}-\bar{u}}$ 无定义；而 CJ15 未拟合奇异压低因子。此处所有矩的不确定度要么基于 68\% C.L.，要么已相应重标定以便比较。
}
\label{tab:moments}
\end{table}

\begin{figure*}
\includegraphics[scale=0.29]{images/mom-comps_v3.pdf}
\includegraphics[scale=0.29]{images/mom-comps_FS_v2.pdf}
\caption{
	表~\ref{tab:moments} 所汇总 PDF 矩的图形比较，但不包括与 Gottfried 求和规则研究相关的 $\bar{d}-\bar{u}$ 组合零阶矩；后者示于
        图~\ref{fig:moments2}。CJ15 全球拟合不从数据独立确定 $\kappa_s$ 与 $\langle x \rangle_{s}$，故未画出其 respective 预言。
}
\label{fig:moments1}
\end{figure*}


\begin{figure*}
\includegraphics[scale=0.45]{images/Gott-comps_v3.pdf}
	\vspace{-1.1cm}
\caption{
	$\int_0^1 dx [\bar{d}(x)-\bar{u}(x)]$ 结果的直观比较。水平的黑色嵌套带与红色带分别对应 1991 年~\cite{Amaudruz:1991at} 与 1994 年~\cite{Arneodo:1994sh} 原 NMC 分析提取的数值。这些提取基于在 $Q^2 = 4$ GeV$^2$、$x \lesssim 0.7$ 范围内测得的氘核-质子结构函数比，直接用夸克-部分子模型抽取味不对称 PDF。虽然该数据集中除最高 $x$ 分箱外都与 CTEQ 运动学割断相容，但其极低的 $Q$ 恰好位于 $Q$ 割断边界，很可能受可观的高阶扭度修正影响，尤其对较高 $x$ 分箱。
}
\label{fig:moments2}
\end{figure*}

我们用更新后的 CT18 与 CT18Z NNLO 拟合计算了一批典型的基准 PDF Mellin 矩，并与较旧的 CT14$_\mathrm{HERAII}$ NNLO 参数化以及近期 MMHT14 NNLO、CJ15 NLO、NNPDF3.1 全球分析比较。
所有情形中，矩均在 $\overline{\mathrm{MS}}$ 因子化标度 $\Q = 2$ GeV 估算，这也是格点 QCD 计算的标准匹配标度。我们把 PDF 矩计算的数值结果汇总在表~\ref{tab:moments} 各条目以及图s.~\ref{fig:moments1}--\ref{fig:moments2} 中。
需要指出，CJ15 NLO 不确定度之所以相对较小，主要归因于使用
$\Delta \chi^2 = 1$ 判据，某些情况下还由于其较为受限的参数化。


{\it 观察。} 总体而言，轻味分布的矩——包括同位矢量（即 $u\!-\!d$）组合的
$\langle x^{1,3} \rangle_{u^+-d^+}$ 与 $\langle x^2 \rangle_{u^--d^-}$——彼此相符。值得注意的是，CT 给出的同位矢量一阶矩 $\langle x \rangle_{u^+-d^+}\! \sim\! 0.156\!-\!0.159$ 略大于本文考虑的其他拟合所得的 $\langle x \rangle_{u^+-d^+}\! \sim\! 0.151\!-\!0.152$，但在 $1\sigma$ 水平上高度吻合。类似地，胶子 PDF 一阶矩的重现一致性也很稳固：可理解为其携带质子在 $\Q = 2$ GeV 标度纵动量的 $\sim\! 41\%$。CT18Z NNLO 下胶子对动量求和规则的总贡献略小，给出 $\langle x \rangle_{g} = 0.407(8)$，但仍轻易地在误差范围内与 CT18 NNLO 的 $\langle x \rangle_{g} = 0.414(8)$ 相符。这与图~\ref{fig:PDFbands1} 右下面板所示 CT18Z 中心胶子的适度减低一致。

对分味夸克密度各自贡献给质子纵动量的份额，我们的新全球分析也与先前和其他拟合的结果普遍强烈趋同。这对 $u^+$ 与 $d^+$ 总一阶矩尤其成立：$\langle x \rangle_{u^+}\! \sim\! 0.35$、$\langle x \rangle_{d^+}\! \sim\! 0.193\!-\!194$。核子总奇异性动量的情况类似，只是围绕 $\langle x \rangle_{s^+}\! =\! 3-4\%$ 的定量散布稍大。CT18 NNLO 给出 $\langle x \rangle_{s^+} = 3.3 \pm 0.9\%$ ——与 CT14$_\mathrm{HERAII}$ 相近。
转到 CT18Z 时，偏好的核子奇异性含量增大 $24\%$，误差相当。


除夸克与胶子密度的一阶矩 $\langle x \rangle_{q^+,g}$ 外，我们还按式~(\ref{eq:qmom}) 估算了若干 $q-\bar{q}$ 夸克不对称的二阶矩，发现 $\langle x^2 \rangle_{u^-,d^-}$ 在 CT18(Z) 与先前计算之间非常接近。近来的 CT 拟合与 CJ15 都不独立参数化 $s$ {\it vs}.~$\overline{s}$，故表~\ref{tab:moments} 省去 $\langle x^2 \rangle_{s^-}$ 条目。


积分 PDF 矩的结果也与唯象求和规则相关——例如 Gottfried 求和
规则~\cite{Gottfried:1967kk}，它把结构函数差 $F^{p-n}_2 = F^{p}_2\! -\! F^{n}_2$ 的 $x^{-1}$ 加权矩
\begin{equation}
\int_0^1\frac{dx}{x}F^{p-n}_2(x,\Q)\big|_\mathrm{QPM} =\frac{1}{3} - \frac{2}{3}\int_{0}^{1} dx\, [\bar{d}-\bar{u}](x,\Q)\ ,
\label{eq:Gott}
\end{equation}
通过 $\mathrm{SU}(2)$ 关系 $\overline{d} = \overline{u}$ 的破缺与轻夸克海中的味对称破缺联系起来。
对与 Gottfried 求和规则相关的零阶矩，我们在 CT18 中得到 $\langle 1 \rangle_{\bar{d}-\bar{u}}\! =\! -0.12\!\pm\!0.35$
（CT18Z 为 $\langle 1 \rangle_{\bar{d}-\bar{u}}\! =\! -0.07\!\pm\!0.29$），大体与其他 PDF 分析一致。其他分析对 $\langle 1 \rangle_{\bar{d}-\bar{u}}$ 给出的不确定度更窄，但这源自其在低 $x$ 区 $x\! \le\! 0.001$ 相对更受限的参数化。该零阶矩由 $\bar{u},\,\bar{d}$ PDF 的低 $x$ 行为主导，而高能数据在那里仍相对稀疏，如图~\ref{fig:ct18data_xmu} 所示。相反，NNPDF 对 $x \to 0$ 时 $\bar{u}$ 与 $\bar{d}$ 的相对行为不加任何限制，因而
$\langle 1 \rangle_{\bar{d}-\bar{u}}$ 无数值定义；表~\ref{tab:moments}
与图~\ref{fig:moments2} 中相应的 NNPDF3.1 条目因此留空。
CT 对轻夸克海使用明显更灵活的参数化（组合的 $\bar u$ 与 $\bar d$ PDF 共 11 个参数，
相比之下 MMHT14 对 $\bar{d}-\bar{u}$ 用 5 参数~\cite{Harland-Lang:2014zoa}，CJ15 对 $\bar{d}/\bar{u}$ 用 5 参数~\cite{Accardi:2016qay}），且不对 $\bar{d}\!-\!\bar{u}$ 的符号施加约束，这可从图~\ref{fig:DBandSBbands} 左上面板的 $\bar{d}/\bar{u}(x,Q\!=\! 1.4\,\mathrm{GeV})$ 比值图看出。因此我们发现，即便具备对 $\langle 1 \rangle_{\bar{d}-\bar{u}}$ 很重要的低 $x$ 区灵活性，现代高能数据仍允许零阶矩取很宽的范围。


CT18(Z) 的 $\langle 1 \rangle_{\bar{d}-\bar{u}}$ 数值与 1991 和 1994 年原 NMC 分析计算的矩~\cite{Amaudruz:1991at,Arneodo:1994sh} 相符；在图~\ref{fig:moments2} 中我们分别以内黑带与外红带表示 1991 年~\cite{Amaudruz:1991at} 与 1994 年~\cite{Arneodo:1994sh} 的提取值。原 NMC 分析的几个方面可以预期会低估 Gottfried 矩的全部实验不确定度，但最主要的是 NMC 仅对区域 $0.004\! <\! x\! <\! 0.8$ 敏感。此外，像式~(\ref{eq:Gott}) 那样直接把 NMC 结构函数矩匹配到 $\langle 1 \rangle_{\bar{d}-\bar{u}}$ 涉及领头阶夸克-部分子计算，必然引入缺失高阶及其他 QCD 效应的修正——它们不在图~\ref{fig:moments2} 的带内。再者，在由氘核-质子截面比确定同位矢量结构函数 $F^{p-n}_2$ 时，NMC 采用相当受限的参数化来执行低 $x$ 外推并表示氘核绝对结构函数。
为便于比较，考察 NMC 实测区域内的 Gottfried 求和规则很有启发；我们发现 CT 与 NMC 合理相符：
\begin{align}
	\frac{1}{3} - \frac{2}{3}\int_{0.004}^{0.8} dx\, [\bar{d}-\bar{u}](x,\Q=2\,\mathrm{GeV}) &= 0.227 \pm 0.016\ (\mathrm{NMC\, '91 }) \\
											       &= 0.221 \pm 0.021\ (\mathrm{NMC\, '94 }) \\
											       &= 0.260 \pm 0.053\ (\mathrm{CT18})\ .
\end{align}
我们强调，就受限积分 $0.004\! <\! x\! <\! 0.8$ 而言 CT 关于 NMC 提取值的不确定度相对较窄，凸显了在探索尚浅的 $x\! <\! 10^{-3}$ 区域内 $\bar{u}, \bar{d}$ 低 $x$ PDF 不确定度的重要性。必须先把它们进一步控制在手，才能让高能数据的唯象分析就式~(\ref{eq:Gott}) 矩水平上的 $\mathrm{SU}(2)$ 味对称破缺作出定论。

\begin{table}
\begin{tabular*}{\textwidth}{c| @{\extracolsep{\fill}} cc|c||c}
\hline
PDF moment                             &   {\bf CT18}   &    {\bf CT18Z}   &    {\bf \CTHERAII}   &   {\bf Lattice}     \tabularnewline
\hline
                                       &                &                  &                   &   $0.153\!-\!0.194^{N_f=2+1+1} \star$             \tabularnewline
$\langle x \rangle_{u^+-d^+}$          & $0.156(7)$     &    $0.156(6)$    &    $0.159(6)$     &   $0.111\!-\!0.209^{N_f=2+1} \star$        \tabularnewline
                                       &                &                  &                   &   $0.166\!-\!0.212^{N_f=2} \star$         \tabularnewline
\hline
$\langle x^2 \rangle_{u^--d^-}$        & $0.055(2)$     &    $0.055(2)$    &    $0.055(2)$     &   $0.107(98) \star$        \tabularnewline
$\langle x^3 \rangle_{u^+-d^+}$        & $0.022(1)$     &    $0.022(1)$    &    $0.022(1)$     &     {\it N/A}             \tabularnewline
\hline
                                       &                &                  &                   &   $0.427(92)$ \cite{Alexandrou:2020sml}        \tabularnewline
$\langle x \rangle_{g}$                & $0.414(8)$     &    $0.407(8)$    &    $0.415(8)$     &   $0.482(69)(48)$ \cite{Yang:2018nqn}        \tabularnewline
                                       &                &                  &                   &   $0.47(4)(11)$ \cite{Yang:2018bft}        \tabularnewline
\hline
\raisebox{-0.25cm}{$\langle x \rangle_{u^+}$}              & \raisebox{-0.25cm}{$0.350(5)$}     &    \raisebox{-0.25cm}{$0.350(4)$}    &    \raisebox{-0.25cm}{$0.351(5)$}     &   \raisebox{-0.12cm}{$0.359(30)$ \cite{Alexandrou:2020sml}}       \tabularnewline
                                       &                &                  &                   &   \raisebox{0.12cm}{$0.307(30)(18)$ \cite{Yang:2018nqn}}        \tabularnewline
\hline
\raisebox{-0.25cm}{$\langle x \rangle_{d^+}$}              & \raisebox{-0.25cm}{$0.193(5)$}     &    \raisebox{-0.25cm}{$0.194(5)$}    &    \raisebox{-0.25cm}{$0.193(6)$}     &   \raisebox{-0.12cm}{$0.188(19)$ \cite{Alexandrou:2020sml}}        \tabularnewline
                                       &                &                  &                   &   \raisebox{0.12cm}{$0.160(27)(40)$ \cite{Yang:2018nqn}}        \tabularnewline
\hline
\raisebox{-0.25cm}{$\langle x \rangle_{s^+}$}              & \raisebox{-0.25cm}{$0.033(9)$}     &    \raisebox{-0.25cm}{$0.041(8)$}    &    \raisebox{-0.25cm}{$0.031(8)$}     &   \raisebox{-0.12cm}{$0.052(12)$ \cite{Alexandrou:2020sml}}        \tabularnewline
                                       &                &                  &                   &   \raisebox{0.12cm}{$0.051(26)(5)$ \cite{Yang:2018nqn}}        \tabularnewline
\hline
$\langle x^2 \rangle_{u^-}$            & $0.085(1)$     &    $0.084(1)$    &    $0.085(1)$     &   $0.117(18)$ \cite{Deka:2008xr}        \tabularnewline
$\langle x^2 \rangle_{d^-}$            & $0.030(1)$     &    $0.029(1)$    &    $0.030(1)$     &   $0.052(9)$ \cite{Deka:2008xr}        \tabularnewline
$\langle x^2 \rangle_{s^-}$            &   ---          &      ---         &      ---          &      {\it N/A}            \tabularnewline
\hline
$\langle 1 \rangle_{\bar{d}-\bar{u}}$  & $-0.12(35)$    &    $-0.07(29)$   &    $-0.37(41)$    &     ---             \tabularnewline
$\kappa_s$                             & $0.49(16)$     &    $0.61(14)$    &    $0.46(13)$     &   $0.795(95)$ \cite{Liang:2019xdx}       \tabularnewline

\end{tabular*}\caption{
	同表~\ref{tab:moments}，但现在把 CT18(Z) 与 CT14$_\mathrm{HERAII}$
	的最新结果与最右列精选的近期格点 QCD 计算比较。后者一般取自文献~\cite{Lin:2017snn,Lin:2020rut} 的近期白皮书。本表信息并非穷举而是摘要；欲广泛了解现代格点计算的读者请参阅文献~\cite{Lin:2017snn,Lin:2020rut} 的详细综述。来自单一计算的格点条目附有相应参考文献；由多个格点提取组合而成的条目以``$\star$''标注。特别是 $\langle x \rangle_{u^+-d^+}$，我们依照文献~\cite{Lin:2020rut} 按所用格点作用量的活跃味数 $N_f$ 分组给出各计算的范围。至于上面所示的 $\langle x^2 \rangle_{u^--d^-}$，它是文献~\cite{Gockeler:2004wp} 结果与文献~\cite{Dolgov:2002zm} 报告的两个独立计算的平均。
}
\label{tab:moments2}
\end{table}

我们可以把 $\overline{d} \neq \overline{u}$ 破缺的分析推广到 $\mathrm{SU}(3)$ 扇区：分析式~(\ref{eq:Rs}) 中诸分布一阶矩之比，得到奇异压低因子矩比
\begin{equation}
	\kappa_s(\Q) \equiv \frac{\langle x \rangle_{s^+}}{\langle x \rangle_{\bar{u}} + \langle x \rangle_{\bar{d}}}\ ,
\label{eq:str-supp}
\end{equation}
如图~\ref{fig:moments1} 所示。表~\ref{tab:moments} 最后一行列出上述 PDF 参数化下该量的数值结果（CJ15 除外，它把 $R_s(x,\Q)$ 设为常数，使 $s^+(x,Q)$ 正比于 $\overline{u}(x,Q) + \overline{d}(x,Q)$）。
在不确定度之内，我们计算的矩与传统奇异压低场景 $\kappa_s\! =\! 0.5$ 大体一致。从 CT14$_\mathrm{HERAII}$ 到 CT18，$Q = 2\,\mathrm{GeV}$ 处偏好的中心值适度增强、相应不确定度增大：由 $\kappa_s({\mathrm{CT14}_\mathrm{HERAII}})\! =\! 0.46\! \pm\! 0.13$
  变为 $\kappa_s(\mathrm{CT18})\! =\! 0.49\! \pm\! 0.16$，与 MMHT14 特别接近。纳入 ATLAS $W,\, Z$ 产生数据以及其他导向 CT18Z 的改变，使比值明显升至 $\kappa_s(\mathrm{CT18Z})\! =\! 0.61\! \pm\! 0.14$，且相对 CT18 其不确定度略有收缩，数值更接近 NNPDF3.1。
最近，$\chi$QCD 合作组在文献~\cite{Liang:2019xdx} 报道了 $\kappa_s$ 的首次格点计算：$\kappa_s(\Q = 2\,\mathrm{GeV}) = 0.795 \pm 0.079\, (\mathit{stat}) \pm 0.053\, (\mathit{sys})$。
确实，该结果恰好落在典型唯象拟合偏好值 $\kappa_s \sim 0.5$ 的上缘之外，但与因纳入 ATLAS 7 TeV 单举 $W,Z$ 产生数据而得出的 CT18Z 结果在 $1\sigma$ 水平相符。

这一点以及表~\ref{tab:moments2} 最右列所列其他在 QCD 格点上确定的 PDF 矩条目，历来都有系统性高估唯象提取值的倾向。较新的格点计算在某些情形已开始逼近唯象矩——例如同位矢量 $u\!-\!d$ 矩，或 $u, d$ 夸克与胶子的总动量 $\langle x \rangle_{u^+,d^+,g}$——其格点不确定度之大也足以允许与全球分析相符。

概略地说，格点上 PDF 矩由三点函数与两点函数关联函数之比提取
\cite{Gockeler:1995wg,Lin:2017snn}：
\begin{equation}
	R(t,\tau,\mathbf{p},\hat{O}) = \frac{\sum_{a,b} \Gamma_{b,a} \langle B_a(t,\mathbf{p}) | \hat{O}(\tau) |
	B_b(0,\mathbf{p}) \rangle}{\sum_{a,b} \Gamma_{b,a} \langle B_a(t,\mathbf{p}) | B_b(0,\mathbf{p}) \rangle}\ ,
\end{equation}
其中 $B_{a,b}$ 是重子插值算符，$t$ 是源-汇欧氏时间间隔，$\tau$ 是与式~(\ref{eq:OPE}) 所述算符 $\hat{O}$ 插入相联系的欧氏时间。对核子部分子分布的低阶矩而言，格点输出很大程度上受两方面相互作用的支配：一是关联函数的激发态污染（一般随欧氏时间按 $\sim\! \exp{(-m_i t)}$ 衰减），二是信噪比（按 $S/N\! \sim\! \exp{(-(E_N -[3/2]m_\pi) t)}$ 衰减）。因此在物理 pion 质量（或其手征外推）下做格点计算，恰好在核子激发态贡献相对受到压制的较大格点时间处，信噪比更快恶化。这些格点效应之间（连同其他系统伪影）的微妙关系，使对表~\ref{tab:moments2} 中目前或大或小的 PDF 矩格点结果的直白解读变得复杂。



\clearpage
